% =====================================================================
%  24.2  Counting paths by distance -- claim each path at its top centroid.
%  Exports: sec:c18-count  def:c18-claim  lst:c18-count
%           prob:c18-715c  lst:c18-715c
%  Solved CF 715C "Digit Tree" (2700), live-verified via the problemset API
%  (tags dfs and similar / divide and conquer / dsu / trees).  Official
%  samples read live:
%    S1: n=6, M=7 -> 7  (note enumerates (0,4),(1,2),(1,5),(3,2),(2,5),(5,2),
%        (3,5) forming 14,21,217,91,7,7,917, all multiples of 7).
%    S2: n=5, M=11 -> 8 (six pairs give 33, two give 3333).
%  Discharges the forward-cites from Ch.22/Ch.23/Ch.16 ("path COUNTING /
%  distance stats -> centroid").  CF 161D (pairs at distance exactly K) is
%  OWNED by Ch.13 as the O(nK) subtree DP -- back-cited by name/label only,
%  never re-linked.  Hand-verified here: the distance-exactly-2 count on T is
%  17 (= sum_z C(deg(z),2)); the top centroid 4 claims 9 of them (6 with the
%  centroid itself + 3 pairs among its three neighbours); recursion claims 8.
% =====================================================================
\section{Counting Paths by Distance}\label{sec:c18-count}

\begin{previouslybox}[Ch.~13 counted distance-\(K\) pairs in \(\bigO{nK}\); here \(K\) is huge]
Chapter~13 harvested, for every node, a histogram of downward distances and merged
children, answering ``how many pairs sit at distance exactly \(K\)'' in
\(\bigO{nK}\) time (\(\probref{c10-harvest-solved}\), CF~161D). That is excellent
when \(K\le 500\) and hopeless when \(n=10^5\) and \(K\) can reach \(n\): the DP
budget is a product, not a bound. The centroid claim removes the product. A path
has two endpoints and, in a tree, a unique route; charge the path to the highest
centroid that lies on it. At that centroid the path splits into two independent
arms, so its length is a \emph{sum of two depths from one node} -- count those
sums with a histogram and you never touch a \(K\)-sized table.
\end{previouslybox}

\begin{topicintro}
Fix a centroid \(c\) of a live component. Every node \(x\) has a distance
\(d_c(x)\) from \(c\). A pair \(\{u,v\}\) whose path passes through \(c\) has
length \(d_c(u)+d_c(v)\), and \(c\) is precisely the first centroid on that path
when \(u\) and \(v\) live in \emph{different} child-components of \(c\). So at
\(c\): walk its child arms one at a time, and for each new node query a running
histogram for partners already seen (in earlier arms, or \(c\) itself), then merge
the new arm into the histogram. Because a partner is always from an
\emph{earlier} arm, same-child pairs are never counted at \(c\) -- they wait for a
deeper centroid. The recursion visits each node \(O(\log n)\) times, so a linear
per-level merge gives \(\bigO{n\log n}\) total, and the same skeleton answers
\(\le K\), \(=K\), or (with a different per-node value) far more exotic path
statistics.

\textbf{Prerequisites.} \(\secref{c18-tree}\) -- the decomposition, the deleted
marker, and \(\mathit{find\_centroid}\). Chapter~16 (\chref{c11}) -- the identity
\(\dist(u,v)=\depth(u)+\depth(v)-2\depth(\lca)\), which here collapses to
\(d_c(u)+d_c(v)\) because the path goes through \(c\).

\textbf{After this section you can}
\begin{itemize}[nosep]
  \item count unordered pairs at distance \(=K\) in \(\bigO{n\log n}\) by the
        query-then-merge histogram, and explain why same-arm pairs are excluded; \hfill\lv{4}
  \item state the top-centroid claim: each path is counted once, at its highest
        centroid (\(\defref{c18-claim}\)); \hfill\lv{4}
  \item recognise a path statistic that is a \emph{sum of two per-arm values} and
        route it through the same skeleton; \hfill\lv{5}
  \item solve CF~715C: the ``value'' is a decimal hash, and coprimality with 10
        turns the concatenation into a matched-key count. \hfill\lv{5}
\end{itemize}
\end{topicintro}

% ---------------------------------------------------------------
\subsection{The top-centroid claim and the distance histogram}
% ---------------------------------------------------------------

\begin{definition}{The top centroid of a path}{c18-claim}
For distinct nodes \(u,v\), the \emph{top centroid} of the path \(u\text{--}v\) is
the centroid lying on that path with the smallest \(\mathit{cd}\) (the highest in
the centroid tree). It exists because the first centroid ever selected inside a
component containing both \(u\) and \(v\) separates them into different child
components, and therefore sits on their path. We say that centroid \emph{claims}
the pair.
\end{definition}

\begin{insightbox}[Why the correction is free]
The textbook way is ``count every pair with \(d_c\) sum \(=K\), then subtract the
pairs that stayed inside one child component.'' The query-then-merge loop
(\(\coderef{c18-count}\)) does the same thing in one pass: a node only ever meets
partners from \emph{earlier} arms and from the centroid, so a same-arm pair is
never formed at \(c\) and there is nothing to subtract after the fact. The
subtraction and the incremental merge are the same idea; the incremental form is
harder to get wrong, so the listing uses it -- and \(\secref{c18-beyond}\)'s map
method reverts to explicit subtraction when the value is not a small integer
array.
\end{insightbox}

\begin{example}{Counting distance-2 pairs on \(T\)}
On the running tree, \(\sum_z\binom{\deg(z)}{2}=17\) pairs sit at distance exactly
\(2\) (any two neighbours of a node are \(2\) apart in a tree, uniquely). Run the
claim at the top centroid \(4\): distances from \(4\) are
\(\{1,2,3\}{=}1\), \(\{5,6,7,8,9,10\}{=}2\), \(\{11,12,13,14\}{=}3\), \(15{=}4\).
A pair through \(4\) with depth-sum \(2\) is either (centroid, a depth-2 node) or
(two depth-1 nodes): \(1\cdot 6 + \binom{3}{2}=6+3=9\). Those \(9\) are
\((4,5),(4,6),(4,7),(4,8),(4,9),(4,10),(1,2),(1,3),(2,3)\). Every other distance-2
pair lives inside one of \(4\)'s arms (\((5,6)\) under~\(1\), \((9,10)\) under~\(3\),
\((10,15)\) under~\(14\), and so on) and is claimed deeper. \(9+8=17\): the pieces
sum to the ground truth. \(\checkmark\)
\end{example}

\begin{lstlisting}[caption={Count unordered pairs at distance exactly K over a fixed tree by centroid decomposition: at each centroid, query-then-merge a depth histogram across the child arms so same-child pairs are never formed here. Reuses the deleted marker and find-centroid from lst:c18-build. ASCII only}, label={lst:c18-count}]
// K given; H[d] = how many seen nodes sit at distance d from the current centroid
const int MAXN = 200005;
vector<int> adj[MAXN];
bool dead[MAXN]; int sz[MAXN];
int K;
long long ans = 0;
int touched[MAXN], tcnt;            // depths written, so we can clear H

int H[MAXN];                        // zero-initialised

void collect_depth(int u, int p, int d, vector<int>& arm){
    arm.push_back(d);
    if(d <= K) touched[tcnt++] = d;
    for(int v : adj[u]) if(v != p && !dead[v]) collect_depth(v, u, d+1, arm);
}
void count_through(int c){
    H[0] = 1;                        // the centroid itself, at distance 0
    for(int v : adj[c]){
        if(dead[v]) continue;
        vector<int> arm;
        collect_depth(v, c, 1, arm);              // depths of this whole arm
        for(int d : arm)
            if(K - d >= 0) ans += H[K-d];         // pair with earlier arms + c
        for(int d : arm) H[d]++;                  // only THEN merge the arm
    }
    H[0] = 0;
    for(int i = 0; i < tcnt; i++) H[touched[i]] = 0;   // clear the histogram
    tcnt = 0;
}
// full recursion (same shape as lst:c18-build):
void solve_comp(int entry){
    int tot = dfs_size(entry, -1);
    int c = find_centroid(entry, -1, tot);
    count_through(c);
    dead[c] = true;
    for(int v : adj[c]) if(!dead[v]) solve_comp(v);
}
// call  solve_comp(1);  then  ans  = #pairs at distance K.
\end{lstlisting}

\begin{complexitybox}[The distance count]
Each centroid does work linear in its component (one depth-collection and one
query per node), and every node lies in \(O(\log n)\) components
(\(\thmref{c18-height}\)), so total time is \(\bigO{n\log n}\) with
\(\bigO{n}\) scratch -- independent of \(K\). That independence is the whole
point versus the \(\bigO{nK}\) subtree DP of Chapter~13: when \(K\) is a
constant \(\le 500\) the DP wins on small \(n\); when \(K\) is \(\Theta(n)\) the
centroid count is the only one that finishes.
\end{complexitybox}

% ---------------------------------------------------------------
\subsection{Problem Lab: CF 715C -- the number read along a path, mod M}
% ---------------------------------------------------------------

\solvedproblem{c18-715c}{Digit Tree}{CF 715C}
  {https://codeforces.com/problemset/problem/715/C}{CF 2700 L5}
  {centroid decomposition; forward/backward decimal hashes; gcd(M,10)=1 gives an inverse of 10; per-arm subtraction over hash maps}

\begin{problemstatement}[Digit Tree (statement and samples verified against the live problem)]
A tree on \(n\) vertices \(0,\dots,n-1\) (\(2\le n\le 10^5\)) has a single nonzero
digit \(1\)–\(9\) on each edge. A positive integer \(M\) (\(1\le M\le 10^9\)) is
given with \(\gcd(M,10)=1\). An \textbf{ordered} pair of distinct vertices
\((u,v)\) is \emph{interesting} if the decimal number obtained by writing the edge
digits along the shortest path from \(u\) to \(v\), \emph{in the order they are
encountered}, is divisible by \(M\). Count the interesting pairs.

\textbf{Examples (official).}
\begin{verbatim}
Input:              Output:
6 7                 7
0 1 2
4 2 4
2 0 1
3 0 9
2 5 7

Input:              Output:
5 11                8
1 2 3
2 0 3
3 0 3
4 3 3
\end{verbatim}
Sample 1's note lists the seven pairs
\((0,4),(1,2),(1,5),(3,2),(2,5),(5,2),(3,5)\) giving \(14,21,217,91,7,7,917\) --
all multiples of \(7\); \((2,5)\) and \((5,2)\) count separately (ordered).
\end{problemstatement}

\begin{attemptgate}[Attempt it yourself first -- \lv{5}, give yourself 45 minutes]
Write the path number as ``\(u\) to \(c\), then \(c\) to \(v\)'' for the top
centroid \(c\) and try to separate the two halves. \peek{the digits from \(u\) up
to \(c\) are the reverse of those from \(c\) down to \(u\); how does appending a
block of digits combine two numbers?}
\nudge{the value read along \(u\to v\) is \(\mathit{back}[u]\cdot 10^{\mathit{dep}[v]} + \mathit{for}[v]\), where \(\mathit{back}[u]\) is the number read \(u\to c\) and \(\mathit{for}[v]\) the number read \(c\to v\). You need that \(\equiv 0 \pmod M\).}
\fullhint{Divide by \(10^{\mathit{dep}[v]}\), legal since \(\gcd(M,10)=1\): the condition
becomes back[u] \(\equiv -\,\text{for}[v]\cdot 10^{-\mathit{dep}[v]}\pmod M\). Store each
node twice -- as a start by key back[u], as an end by key
\(-\,\text{for}[v]\cdot\) inv10\([\mathit{dep}[v]]\) -- and count matched keys with a hash
map. Count all ordered matches over the component, then \emph{subtract} the
matches inside each single child arm (those paths do not pass through \(c\)); the
centroid, with for \(=\) back \(= 0\), pairs correctly with every node whose
one-sided value is \(0\pmod M\). Recurse; total \(\bigO{n\log n}\) expected.}
\end{attemptgate}

\solstep{Understanding}
The unit of work is still a path claimed at its top centroid, but the per-node
value is no longer an integer distance: it is a residue that depends on the
\emph{digits} and their \emph{order}, so a path through \(c\) is a concatenation of
an upward part and a downward part. Concatenation is `\(A\cdot 10^{\mathit{len}(B)}+B\)',
and coprimality with 10 makes \(10^{\mathit{len}(B)}\) invertible, which decouples the two
sides into an equality of single-node keys. That decoupling is exactly what lets
the \(\secref{c18-count}\) skeleton -- count by a matched key, subtract one-arm
pairs -- survive the move from distances to decimal numbers.

\solstep{Key Observations}
\begin{observation}
For a path \(u\to v\) through centroid \(c\), let \(\mathit{back}[u]\) be the value
of the digits read \(u\to c\) and \(\mathit{for}[v]\) the value read \(c\to v\),
with \(\mathit{dep}[\cdot]\) the edge counts. Then
\(\mathrm{val}(u,v)=\mathit{back}[u]\cdot 10^{\mathit{dep}[v]}+\mathit{for}[v]\).
\end{observation}
\begin{observation}
Since \(\gcd(M,10)=1\), \(10^{\mathit{dep}[v]}\) has an inverse mod \(M\), so
\(\mathrm{val}(u,v)\equiv 0 \iff
\mathit{back}[u]\equiv -\,\mathit{for}[v]\cdot 10^{-\mathit{dep}[v]}\pmod M\): a
pure key-equality, with \(u\)'s key independent of \(v\).
\end{observation}
\begin{observation}
A pair is claimed by \(c\) only if its two endpoints are in different child arms
(or one endpoint is \(c\)). Matches with both endpoints in one arm are subtracted;
they are counted at a deeper centroid, exactly as in \(\secref{c18-count}\).
\end{observation}

\solstep{The Approach Ladder}
\begin{center}\renewcommand{\arraystretch}{1.25}\setlength{\tabcolsep}{3pt}
\begin{tabular}{@{}c p{0.28\linewidth} p{0.30\linewidth} l c@{}}
\toprule
\textbf{\#} & \textbf{Approach} & \textbf{Idea} & \textbf{Time} & \textbf{Verdict}\\
\midrule
1 & \lv{2}\ all ordered pairs & DFS a number per \((u,v)\) & \(\bigO{n^2}\) & \tle \\
2 & \lv{4}\ rooted DP on digits & fix a root, fold & -- & \textcolor{warningred}{\ding{55}\ \textbf{WA}} \\
3 & \lv{5}\ \textbf{centroid + hash keys} & back / for keys, subtract arms (\(\coderef{c18-715c}\)) & \(\bigO{n\log n}\) & \ac \\
\bottomrule
\end{tabular}\end{center}
Rung~2 fails because a single root cannot see a path that climbs to the top
centroid and back down with the correct digit order in both directions; the
reverse-direction hash \(\mathit{back}[\cdot]\) is not derivable from
\(\mathit{for}[\cdot]\) without the \(10^{-\mathit{dep}}\) normalisation the centroid
framing supplies.

\solstep{Solution}
Precompute \(10^k\bmod M\) and its inverses. Build the centroid decomposition
(\(\coderef{c18-build}\)); at each centroid \(c\) collect, for every node of the
component, the pair of keys \((\mathit{back},\mathit{key})\); add
\(\mathit{pairs}(\text{all})\) and subtract \(\mathit{pairs}(\text{each arm})\),
where \(\mathit{pairs}\) counts ordered distinct matches of the start-key against
the end-key.

\solstep{Implementation}
\begin{lstlisting}[caption={CF 715C: centroid decomposition with a forward/backward decimal hash per node and a per-arm subtraction over hash maps; gcd(M,10)=1 supplies inv(10, M). M=1 is short-circuited. ASCII only}, label={lst:c18-715c}]
#include <bits/stdc++.h>
using namespace std;
typedef long long ll;
const int MAXN = 100005;
int n; ll M;
vector<pair<int,int>> adj[MAXN];        // (neighbour, digit 1..9)
bool dead[MAXN]; int sz[MAXN];
ll p10[MAXN], ip10[MAXN];               // 10^k mod M, and its inverse
ll ans = 0;

ll egcd(ll a, ll b, ll &x, ll &y){
    if(!b){ x=1; y=0; return a; }
    ll g, x1, y1; g = egcd(b, a % b, x1, y1);
    x = y1; y = x1 - (a / b) * y1; return g;
}
ll invmod(ll a){ ll x, y; egcd(a, M, x, y); return (x % M + M) % M; }
ll norm(ll x){ x %= M; if(x < 0) x += M; return x; }

int dfs_size(int u, int p){
    sz[u] = 1;
    for(auto &e : adj[u]) if(e.first != p && !dead[e.first])
        sz[u] += dfs_size(e.first, u);
    return sz[u];
}
int walk_cent(int u, int p, int tot){
    bool moved = true;
    while(moved){
        moved = false;
        for(auto &e : adj[u]){
            int v = e.first;
            if(v != p && !dead[v] && sz[v] > tot/2){ p = u; u = v; moved = true; break; }
        }
    }
    return u;
}

struct Info { ll b, k; };               // b = back[u]; k = -for*inv10[dep]
void collect(int u, int p, int dep, ll fwd, ll bwd, vector<Info>& out){
    out.push_back({ norm(bwd), norm(-fwd * ip10[dep]) });
    for(auto &e : adj[u]){
        int v = e.first; ll w = e.second;
        if(v == p || dead[v]) continue;
        collect(v, u, dep+1, norm(fwd*10 + w), norm(w * p10[dep] + bwd), out);
    }
}
ll pairs(const vector<Info>& a){         // ordered distinct (u,v), back[u]==key[v]
    unordered_map<ll,ll> fB, fK;
    for(auto &x : a){ fB[x.b]++; fK[x.k]++; }
    ll cnt = 0;
    for(auto &kv : fB){ auto it = fK.find(kv.first); if(it != fK.end()) cnt += kv.second * it->second; }
    for(auto &x : a) if(x.b == x.k) cnt--;        // drop u == v
    return cnt;
}
void process(int c){
    vector<Info> all; all.push_back({ 0, 0 });     // the centroid itself
    vector<vector<Info>> arms;
    for(auto &e : adj[c]){
        int v = e.first; ll w = e.second;
        if(dead[v]) continue;
        vector<Info> arm;
        collect(v, c, 1, norm(w), norm(w), arm);
        all.insert(all.end(), arm.begin(), arm.end());
        arms.push_back(arm);
    }
    ans += pairs(all);
    for(auto &a : arms) ans -= pairs(a);           // subtract same-child pairs
}
void decompose(int entry){
    int tot = dfs_size(entry, -1);
    int c = walk_cent(entry, -1, tot);
    process(c);
    dead[c] = true;
    for(auto &e : adj[c]) if(!dead[e.first]) decompose(e.first);
}
int main(){
    scanf("%d %lld", &n, &M);
    if(M == 1){ printf("%lld\n", (ll)n * (n - 1)); return 0; }   // every pair works
    for(int i = 0; i < n-1; i++){
        int u, v, w; scanf("%d %d %d", &u, &v, &w);
        adj[u].push_back({v, w}); adj[v].push_back({u, w});
    }
    ll inv10 = invmod(10 % M);
    p10[0] = 1 % M; ip10[0] = 1 % M;
    for(int i = 1; i <= n; i++){ p10[i] = p10[i-1]*10 % M; ip10[i] = ip10[i-1]*inv10 % M; }
    decompose(0);
    printf("%lld\n", ans);
}
\end{lstlisting}

\solstep{Trace}
\begin{tracebox}[Sample 2: \(n=5\), \(M=11\), every edge digit \(=3\)]
A run of \(k\) threes is the number \(3,33,333,\dots\); mod \(11\), \(33\equiv 0\)
and \(3333\equiv 0\) while \(3\not\equiv 0\) and \(333\equiv 3\), so a path reads to
\(0\pmod{11}\) exactly when it uses an \emph{even} number of edges. The four edges
\(1\!-\!2,2\!-\!0,3\!-\!0,4\!-\!3\) form the path \(1\!-\!2\!-\!0\!-\!3\!-\!4\)
(order \(1,2,0,3,4\)); a pair is interesting iff their positions differ by an even
number. Position gaps of \(2\) give \(\{1,0\},\{2,3\},\{0,4\}\) and a gap of \(4\)
gives \(\{1,4\}\): four unordered pairs, eight ordered ones --
\((4,0),(0,4),(3,2),(2,3),(0,1),(1,0),(4,1),(1,4)\), the official eight, six reading
\(33\) and two reading \(3333\). The centroid decomposition returns \(\boxed{8}\)
whatever centroid is chosen first, since each pair is claimed at exactly one top
centroid.
\end{tracebox}

\begin{warningbox}[Ordered pairs, and the modulo must be normalised]
This problem counts \emph{ordered} pairs: \((2,5)\) and \((5,2)\) are different, and
they can even carry different numbers (\(7\) vs \(7\) here, but not in general), so
\(\mathit{pairs}\) matches the start-key against the end-key over all nodes in both
directions -- never halve the result as you would for an unordered distance count.
And every key must be pushed into \([0,M)\) with \(\mathit{norm}\): the raw
\(-\,\text{for}\cdot\mathit{inv10}\) is negative, and comparing a negative residue
against a stored positive one silently under-counts -- a classic \textbf{WA} that
passes the small samples and dies on the large ones.
\end{warningbox}

\begin{hardeningbox}
\begin{itemize}[nosep,leftmargin=1.4em]
  \item \textbf{\(\gcd(M,10)=1\) is load-bearing.} The \(10^{-\mathit{dep}}\) normalisation
        needs an inverse of 10; if the input had \(2\mid M\) or \(5\mid M\) the key
        split is invalid. Here it is promised. Guard \(M=1\) separately (every
        ordered pair works: \(n(n-1)\)).
  \item \textbf{Overflow.} \(M\le 10^9\) so products like \(\mathit{fwd}\cdot
        \mathit{inv10}\) reach \(10^{18}\) -- use \code{long long} and take \(\bmod M\)
        immediately; \(\mathit{ans}\) can be \(\sim n^2\), also \code{long long}.
  \item \textbf{Map hygiene.} Building two \code{unordered\_map}s per centroid is
        \(O(\text{component})\); reserve them or use a sorted vector if the judge
        anti-hashes. The subtraction is over the \emph{same} \(\mathit{pairs}\) helper
        -- do not forget to drop \(u=v\) in both the \(\mathit{all}\) and each arm call.
  \item \textbf{Digit vs edge.} The single digit lives on the \emph{edge}, so a path
        of \(t\) edges reads \(t\) digits and has length \(t\); \(\mathit{dep}\) counts edges,
        matching the exponent of 10 exactly.
\end{itemize}
\end{hardeningbox}

\begin{recallbox}[24.2 in one breath]
\emph{Claim every path at its top centroid (\(\defref{c18-claim}\)); there its length
is a sum of two per-arm depths, counted by a query-then-merge histogram in
\(\bigO{n\log n}\) (\(\coderef{c18-count}\)) -- the upgrade of Ch.~13's \(\bigO{nK}\)
DP for CF~161D. The same skeleton survives any path value that factors into two
independent per-arm pieces: CF~715C reads the path as a decimal number, and
coprimality with 10 turns `concatenate and test divisible' into `match a forward
hash against a normalised backward hash, then subtract the same-child pairs'
(\(\coderef{c18-715c}\)).}
\selftest{
  \item Why does query-then-merge across arms already exclude same-child pairs? \emph{(A node only meets partners from strictly earlier arms and the centroid, so an intra-arm pair is never formed here (\(\secref{c18-count}\)).)}
  \item In CF~715C what makes the concatenation separable into single-node keys? \emph{\(\gcd(M,10)=1\), giving \(10^{-\mathit{dep}[v]}\) to move the length factor onto one side (\(\coderef{c18-715c}\)).)}
  \item Which pairs does the per-arm subtraction remove? \emph{Those with both endpoints in one child component of \(c\); they are claimed at a deeper centroid.}
}
\end{recallbox}
